\documentclass[10pt]{article}
\usepackage[margin=0.70in]{geometry}
\usepackage[T1]{fontenc}
\usepackage[utf8]{inputenc}
\usepackage[hidelinks]{hyperref}
\setlength{\parindent}{0pt}
\setlength{\parskip}{0.38em}
\setlength{\emergencystretch}{2em}

\newcommand{\questionsep}{\par\vspace{0.75em}}

\begin{document}

\begin{center}
{\Large Response to the Area Chair}\\[2pt]
Submission 13175
\end{center}

We thank the Area Chair and reviewers.
\textbf{sMMC-22M advances spatial transcriptomics from spot averages toward
cell-centered multimodal analysis by linking morphology, RNA, location, and
biological context at scale.}

Below, we summarize each central question, our response, and the resulting
evidence.

\section*{Reviewer 7Gar: contribution, experimental purpose, and validation}

\begin{itemize}
\item \textbf{Question---What is genuinely new? Response---}
We completed a mutually exclusive provenance audit. The release contains
\textbf{28,315,247 aligned targets from 99 unique samples}, including
\textbf{53,989 newly measured native DBiC-seq cells from 21 in-house samples}.
Only 32 samples overlap HEST-1K and two additional samples overlap
STimage-1K4M only; approximately \textbf{70\% is absent from HEST-1K}. This
separates new measurements, public aggregation, and processed derivatives.

\item \textbf{Question---What do the three experimental axes demonstrate?
Response---}
\begin{itemize}
\item \textit{Scale} tests whether more paired data improve foundation
encoders. Across seven pathology foundation encoders, using 5\%, 10\%, 25\%,
and 100\% of the data increases mean Gene Pearson from
\textbf{0.177 to 0.205, 0.241, and 0.278}. This validates the value of scale.
\item \textit{Resolution} tests whether cell alignment preserves variation
lost through regional averaging.
\item \textit{Rich context} tests whether metadata reveal failures hidden by
average scores. Across assay, dataset, and site contexts,
\textbf{Average--Worst balanced-accuracy gaps reach 0.258--0.975}.
\end{itemize}
Each axis now has an explicit question, design, metric, and conclusion.

\item \textbf{Question---How are derived profiles and textual records
validated? Response---}
Visium HD targets use registered native \(2\,\mu\mathrm{m}\) bins intersected
with cell masks; conflicts are resolved deterministically and unsupported
cells removed. Across \textbf{9,000 polygons} from lung, brain, and pancreas,
\(\pm1\,\mu\mathrm{m}\) perturbations give bin Jaccard
\textbf{0.706--0.816} and expression cosine \textbf{0.936--0.994}. GPT-4o
converts supplied metadata into readable descriptions; PLIP/CONCH audits give
AUC \textbf{0.993/0.988} and \textbf{100\% Top-3 retrieval}. Full prompts,
provenance, QC, and examples will be added to the Appendix.
\end{itemize}

\section*{Reviewer kUmZ: target validity, benchmark breadth, and biological
utility}

\begin{itemize}
\item \textbf{Question---Which targets are natively cell resolved? Response---}
We separate \textbf{16,314,129 native Xenium cells} and \textbf{53,989 native
DBiC-seq cells} from \textbf{7,629,697 derived Visium HD cell-aligned targets},
using \textit{native cell-resolved} and \textit{derived cell-aligned}.

\item \textbf{Question---Does evaluation support the resource's breadth?
Response---}
We now report complete benchmark results for \textbf{all 25 organ categories}
under a unified evaluation protocol. We additionally provide \textbf{eight
cross-sample transfer results} in which the complete test sample is unseen.
Together, these experiments replace the original small set of representative
cases with release-wide and cross-sample evidence.

\item \textbf{Question---What does cell alignment add beyond spot evaluation?
Response---}
Across six native-Xenium samples, cell/8/16/\(55\,\mu\mathrm{m}\) targets
yield Gene Pearson \textbf{0.365/0.365/0.363/0.330}. Yet
\(55\,\mu\mathrm{m}\) pseudo-spots mix cell types in \textbf{55.5--66.4\%}
of dense-tissue regions and affect \textbf{73.8--81.0\%} of cells.
\textbf{Cell alignment preserves cellular identity and local heterogeneity
obscured by regional averaging.}
\end{itemize}

\section*{Reviewers vCPF and aSvP}

We appreciate their constructive suggestions. vCPF requested HD boundary
validation, broader organ coverage, and spatial/metadata calibration; these
are addressed by the audits, 25-category benchmark, and added mean,
coordinate, spatial-KNN, and segmentation-metadata baselines. aSvP asked how
scale and heterogeneity translate into research value; the scaling,
release-wide, cell-versus-spot, and context-gap analyses provide this evidence.

\section*{Summary}

\textbf{sMMC-22M provides an audited, release-wide foundation connecting
spatial transcriptomics with cell-centered computational biology.} Moving
from spots toward cells enables cell-type/state prediction, cell--cell
communication, morphology--RNA analysis, perturbation-response modeling, and
spatially grounded virtual-cell development---questions difficult when
heterogeneous cells are averaged into one spot.

Following the reviewers' suggestions, we will clarify STBoost and the
experimental objectives; distinguish native, processed, and generated
records; add the release-wide results, audits, fixed splits, QC, prompts, and
examples; and release the code and in-house data. \textbf{The new evidence
answers the central questions about novelty, target validity, breadth,
heterogeneity, context, and reproducibility. We hope it resolves the reviewers'
concerns and demonstrates the utility of sMMC-22M.} We thank the Area Chair
and reviewers for recognizing its value and helping us strengthen it.

\end{document}
